\section{Diseño del Control de Seguridad (Nivel 0)}

El Nivel 0 representa la capa de protección más baja, pero con mayor jerarquía dentro de la
arquitectura de control. Su función exclusiva es salvaguardar la integridad de la grúa, la carga
y el personal. Este nivel tiene la autoridad para sobrescribir cualquier comando proveniente del
Nivel 1, cortando la energía de los actuadores y aplicando los frenos mecánicos ante cualquier
anomalía.

A diferencia de los niveles superiores, el control de seguridad opera de forma permanente y
supervisa la planta mediante la redundancia de sensores (finales de carrera de emergencia,
detectores de sobrevelocidad, botón físico de emergencia y el latido de comunicación). La lógica
se modeló mediante dos procesos paralelos: el autómata de estados de seguridad, que decide qué
frenos aplicar ante cada falla, y el monitor de comunicaciones (watchdog), que vigila el enlace
con el control supervisor.

\subsection{Autómata de Estados de Seguridad}

La máquina de estados \texttt{ESTADOS\_SEGURIDAD} supervisa constantemente el estado físico de la
grúa. Durante la operación normal permanece en el estado \texttt{SISTEMA\_OK}, con ambos frenos
energizados (abiertos): el de operación del carro (\textcolor{signalcolor}{$BRK_t$} $= 1$) y el de
emergencia de izaje (\textcolor{signalcolor}{$BRK_{hE}$} $= 1$).

Dentro de este estado nominal, el bloque realiza además un cálculo de seguridad redundante:
recomputa en tiempo real el límite de velocidad de izaje permitido (\textcolor{signalcolor}{vh\_lim})
en función de la masa de la carga, según la curva de potencia constante del motor. De esta forma,
la restricción de sobrevelocidad por debilitamiento de campo queda garantizada de manera
independiente al cálculo que realiza el Nivel 1 (defensa en profundidad).

Ante la detección de una condición crítica, el autómata transiciona a un estado de emergencia
específico, aplicando inmediatamente los frenos correspondientes. Una característica de diseño es
el \textit{frenado selectivo}: solo se detiene el eje afectado por la falla, dejando libre el otro
para no introducir esfuerzos innecesarios. Las condiciones de falla contempladas son:

\begin{itemize}
    \item \textbf{Emergencia de izaje:} se activa uno de tres estados según corresponda
    (\texttt{EMERGENCIA\_IZAJE\_SUP}, \texttt{EMERGENCIA\_IZAJE\_INF} o \texttt{EMERGENCIA\_IZAJE\_SV}):
    si se accionan los finales de carrera de emergencia superior o inferior del izaje, o si la
    velocidad observada supera en un 15\,\% al límite dinámico calculado (sobrevelocidad). En todos
    los casos se cierra el freno de emergencia del izaje (\textcolor{signalcolor}{$BRK_{hE}$} $= 0$),
    manteniendo el carro libre (\textcolor{signalcolor}{$BRK_t$} $= 1$).

    \item \textbf{Emergencia de carro:} se activa uno de tres estados (\texttt{EMERGENCIA\_CARRO\_DER},
    \texttt{EMERGENCIA\_CARRO\_IZQ} o \texttt{EMERGENCIA\_CARRO\_SV}): si el carro acciona los finales
    de carrera de emergencia izquierdo o derecho, o si supera el umbral de sobrevelocidad. Se cierra
    el freno de operación del carro (\textcolor{signalcolor}{$BRK_t$} $= 0$), dejando libre el izaje
    (\textcolor{signalcolor}{$BRK_{hE}$} $= 1$).

    \item \textbf{Emergencia manual:} detención forzada accionada por el operador mediante el botón
    físico de emergencia. En este caso se cierran ambos frenos
    (\textcolor{signalcolor}{$BRK_t$} $= 0$, \textcolor{signalcolor}{$BRK_{hE}$} $= 0$), llevando la
    grúa a una parada total.

    \item \textbf{Emergencia de comunicación:} disparada por la pérdida de comunicación con el
    control supervisor (\textcolor{signalcolor}{FALLA\_WD}), detectada por el proceso watchdog.
\end{itemize}

Una vez mitigada la causa de la parada, el sistema exige un rearme manual mediante el botón de
reset, que hace transicionar la máquina al estado \texttt{MODO\_RECUPERACION} (con los frenos
nuevamente energizados) antes de permitir el regreso a la operación nominal.

\begin{figure}[H]
    \centering
    \includegraphics[width=0.9\textwidth]{nivel_0/estados_seguridad.png}
    \caption{Máquina de estados de seguridad en Stateflow}
    \label{fig:n0_estados_seguridad}
\end{figure}

\subsection{Monitor de Comunicaciones (Watchdog)}

Este segundo proceso paralelo funciona como un supervisor de la integridad del enlace de datos
con el Nivel 1. El control supervisor debe emitir una señal  \textit{heartbeat} que
alterna su estado constantemente.

La máquina de estados \texttt{WATCHDOG} alterna entre \texttt{HEARTBEAT\_DOWN} y
\texttt{HEARTBEAT\_UP} persiguiendo dicha señal, reiniciando un contador de ciclos
(\textcolor{signalcolor}{WD\_counter}) en cada conmutación válida. Si la señal se congela y el
sistema permanece en un solo estado durante una cantidad de ciclos que supere el límite seguro
(\textcolor{signalcolor}{counter\_limit}), se transiciona al estado \texttt{DETECCION\_FALLO} y se
declara una falla de red (\textcolor{signalcolor}{FALLA\_WD} $= 1$). Esta bandera dispara la parada
de emergencia en el autómata principal, evitando que la grúa continúe moviéndose con el último
comando recibido ante una pérdida de comunicación.

\begin{figure}[H]
    \centering
    \includegraphics[width=0.7\textwidth]{nivel_0/watchdog.png}
    \caption{Monitor de comunicaciones (watchdog) en Stateflow}
    \label{fig:n0_watchdog}
\end{figure}

\subsection{Modelo del Control de Seguridad en Simulink}

La lógica de seguridad se concentra en un único bloque Chart (\texttt{Control Nivel 0}). Sus
puertos de entrada reciben el botón de emergencia y de reset, los finales de carrera de
emergencia de ambos ejes, la velocidad observada del izaje, la masa estimada del contenedor, la
señal de recuperación del supervisor y el latido de comunicación. Como salidas entrega los
comandos de freno \textcolor{signalcolor}{$BRK_{hE}$} y \textcolor{signalcolor}{$BRK_t$} (que
sobrescriben a los del Nivel 1) y el conjunto de banderas de falla que informan al supervisor y a
la interfaz el motivo de la parada.

\begin{figure}[H]
    \centering
    \includegraphics[width=0.9\textwidth]{nivel_0/modelo_seguridad.png}
    \caption{Modelo del control de seguridad (Nivel 0) en Simulink}
    \label{fig:n0_modelo}
\end{figure}
